\section{Central Thesis: A Deterministic Diagram Compiler}
\label{sec:central-thesis}

This document specifies a \emph{deterministic, themeable, agent-authorable diagram compiler}: a system that transforms small declarative intermediate representations into presentation-quality visual output through a layered architecture of grammars, a shared kernel, and pluggable rendering backends.

\subsection{The Engine-as-Asset Insight}

The strategic insight underlying this specification is a reframe: \textbf{the real asset is not ``a timeline tool.''} It is a \emph{compiler architecture} capable of hosting a family of visual grammars, each small enough to be reliably emitted by an LLM, each compiling down to a shared Scene IR that guarantees determinism and enables theming across all diagram types.

Timeline is \emph{grammar \#1}---the first proof of the engine, not the whole product. The architecture generalizes to:

\begin{itemize}
  \item \textbf{Flow/pipeline diagrams}: ordered icon+label nodes with arrows, loops, and legends
  \item \textbf{Node-link graphs}: trees, networks, hub-and-spoke, architecture diagrams
  \item \textbf{Comparison matrices}: N columns with checkmarks, stats, or feature rows
  \item \textbf{Stat callouts}: big number + caption compositions
  \item \textbf{Step-card sequences}: numbered ordinal + icon + bullet lists
  \item \textbf{Multi-panel posters}: compositions that arrange multiple grammars into a single visual
\end{itemize}

The value proposition is the \emph{compounding payoff}: animation, themes, new backends, and quality-gates added once in the kernel are inherited by every grammar. A new grammar (e.g., flow diagrams) gains determinism, theming, SVG/PNG/PDF output, and validate-before-render for free.

\subsection{The Pipeline Model}

Every diagram in the compiler follows the same pipeline:

\begin{figure}[h]
\centering
\begin{Verbatim}[frame=single,fontsize=\small]
  Domain IR (grammar-specific)
       |
       v  [Grammar-specific Layout Engine]
       |
  Scene IR (universal primitives)
       |
       v  [Rendering Backend]
       |
  Output (SVG / PNG / PDF / PPTX)
\end{Verbatim}
\caption{The universal pipeline: Domain IR → Layout Engine → Scene IR → Backend → Output}
\label{fig:central-pipeline}
\end{figure}

\paragraph{Validate-before-render.}
The pipeline enforces a strict contract: no diagram reaches the renderer until it passes schema validation and semantic checks. This is not defensive coding; it is a \emph{specification-level guarantee}. An invalid IR produces a diagnostic, never a corrupt visual.

\paragraph{Byte-identical golden determinism.}
The same IR plus the same theme produces byte-identical output across runs, platforms, and time. This property enables CI/CD integration, golden-image testing, and reliable agent round-trips.

\subsection{Why ``Diagram Compiler,'' Not ``Diagram Tool''}

The term ``compiler'' is precise and intentional:

\begin{description}
  \item[Structured input, structured output] The system does not accept freeform drawings or WYSIWYG manipulation. It accepts a declarative specification (the IR) and produces a deterministic artifact.
  
  \item[Separation of concerns] Content (IR), styling (theme), and rendering (backend) are cleanly separated. The IR describes \emph{what} to communicate; the theme describes \emph{how} it looks; the backend describes \emph{where} it goes.
  
  \item[Typed intermediate representations] Like a programming-language compiler, this system uses typed IRs at multiple levels (Domain IR, Scene IR) with well-defined transformations between them.
  
  \item[Determinism by construction] The compiler is a pure function from (IR, theme, backend) to output. No random state, no user interaction, no system clock.
\end{description}

This is not a design tool (Figma, Canva), not a charting library (D3, Plotly), and not a project-management tool (Jira, MS Project). It is a \emph{rendering compiler} for technical explainer diagrams.

\subsection{The Target Niche: Technical Explainer Content}

The inspiration corpus (Section~\ref{sec:corpus-taxonomy}) reveals a coherent niche: \emph{AI-engineering / developer technical-explainer infographics}. These artifacts share visual vocabulary (flows, trees, matrices, step-cards, stat callouts) and share an audience (technical practitioners, executives receiving technical briefings).

This niche is defensible:

\begin{itemize}
  \item \textbf{Small grammars are the feature.} A narrow grammar is both shippable and reliably LLM-emittable. A god-schema that tries to express everything is unmaintainable and hostile to agent generation.
  
  \item \textbf{Automatic layout is required.} The author (human or agent) should not hand-place nodes. The grammar must be constrained enough that a deterministic layout algorithm can position elements without ambiguity.
  
  \item \textbf{The author is often an LLM.} Every design decision must ask: can an agent emit this reliably? If the answer is ``only with special prompting,'' the design is wrong.
\end{itemize}

The explicit \emph{non-goal} is competing with general-purpose freeform canvas tools (Canva, Figma, Excalidraw). Those tools optimize for human creativity and manual placement. This compiler optimizes for deterministic, automated, agent-driven generation of \emph{constrained} visual vocabularies.

\subsection{Elevating the Timeline Grammar Thesis}

Section~\ref{sec:grammar-abstraction-old} (now Part~III) established the ``Timeline Grammar vs.\ Task Scheduling Grammar'' distinction. That thesis remains valid but is now a \emph{special case} of a broader principle:

\begin{quote}
\emph{Every grammar in the compiler adopts a visual-communication abstraction, not an operational/computational abstraction. The IR describes what to show, not what to compute.}
\end{quote}

For timelines, this means: no dependency scheduling, no resource leveling, no critical-path analysis. For flow diagrams, this means: no execution semantics, no data-flow analysis, no simulation. For graphs, this means: no graph algorithms beyond layout---the IR specifies nodes and edges, not traversal results.

The compiler renders \emph{narratives} about technical concepts. It does not execute those concepts.

\subsection{The Kernel/Grammar/Composition Architecture}

The architecture has three layers (detailed in Part~\ref{part:architecture}):

\begin{description}
  \item[Kernel (Part~\ref{part:kernel})] The shared infrastructure: Scene IR, rendering backends, theme system, determinism contract, icon registry, asset loader, lint/quality-gate framework. The kernel knows nothing about timelines, flows, or graphs---it renders primitives.
  
  \item[Grammars (Part~\ref{part:grammars})] Peer modules atop the kernel. Each grammar defines a Domain IR and one or more layout engines that compile Domain IR → Scene IR. Timeline is grammar \#1; future grammars include flow, graph, comparison, stat, step-cards.
  
  \item[Composition (Part~\ref{part:composition})] A thin layer that arranges multiple panels into a poster. Each panel references a grammar's Domain IR. This is how the multi-section infographics are assembled.
\end{description}

\subsection{Summary of the Strategic Evolution}

This specification captures a major strategic evolution:

\begin{enumerate}
  \item \textbf{From ``timeline tool'' to ``diagram compiler.''} The product is the engine, not the first grammar.
  
  \item \textbf{From one IR to two IR layers.} Domain IRs are grammar-specific and small; all compile to the shared Scene IR.
  
  \item \textbf{From SVG-first to Scene-first.} SVG is one backend, not the universal root. The Scene IR is the render contract.
  
  \item \textbf{From timeline-only to grammar family.} Flows, graphs, comparisons, stats, and compositions are in scope.
  
  \item \textbf{From static-only to animation-aware.} Animation is additive, declarative, and backend-conditional.
  
  \item \textbf{From monolith to layered kernel.} The kernel is extractable; grammars are thin consumers.
\end{enumerate}

The following parts specify each layer in detail.
